\documentclass[11pt,a4paper]{article}
\usepackage[margin=20mm,top=22mm,bottom=22mm]{geometry}
\usepackage[spanish,es-nodecimaldot]{babel}
\usepackage{fontspec}
\setmainfont{Calibri}
\setmonofont{Consolas}[Scale=.87]
\usepackage{amsmath,microtype,tabularx,booktabs,array,enumitem,xcolor,hyperref,fancyhdr}
\definecolor{tinta}{HTML}{131B32}
\definecolor{azul}{HTML}{1648FF}
\definecolor{gris}{HTML}{525E77}
\hypersetup{colorlinks=true,linkcolor=azul,urlcolor=azul,pdftitle={Red Pepper | Informe técnico de base de datos},pdfsubject={Normalización y piloto PostgreSQL en Supabase}}
\color{tinta}
\setlength{\parindent}{0pt}
\setlength{\parskip}{6pt}
\setlength{\headheight}{15pt}
\renewcommand{\arraystretch}{1.24}
\setlist[itemize]{leftmargin=1.3em,topsep=1pt,itemsep=2pt,parsep=0pt}
\newcolumntype{Y}{>{\raggedright\arraybackslash}X}
\newcommand{\codigo}[1]{\texttt{#1}}
\newcommand{\seccion}[2]{\vspace{5pt}{\Large\bfseries #1\enspace #2}\par\vspace{2pt}}
\newcommand{\nota}[1]{\par\smallskip{\color{gris}\small #1}\par}
\pagestyle{fancy}
\fancyhf{}
\lhead{\small\color{gris} RED PEPPER / BASE DE DATOS}
\rhead{\small\color{gris} 24 DE NOVIEMBRE DE 2026}
\lfoot{\small\color{gris} Informe técnico · Piloto en Supabase}
\rfoot{\small\color{gris} \thepage}
\renewcommand{\headrulewidth}{.35pt}

\begin{document}

{\color{azul}\bfseries NEXLABS / RED PEPPER}\par
\vspace{6pt}
{\fontsize{34}{37}\selectfont\bfseries Normalización y piloto\\ en PostgreSQL}\par
\vspace{5pt}
{\large Clientes, catálogo y pedidos de Red Pepper}\par
\vspace{5pt}
{\color{gris} Asignatura Base de Datos · Docente Ing. Fabrizio Calienes}

\seccion{1.}{Alcance del modelo}
Red Pepper registra pedidos de clientes y los productos de cada comanda. El esquema \codigo{salchipaperia} en Supabase organiza el catálogo en snacks, bebidas y hamburguesas. Una muestra de 100 pedidos ficticios permite comprobar las relaciones y los importes de la carta.

\begin{tabularx}{\linewidth}{@{}lYY@{}}
\toprule
\textbf{Indicador} & \textbf{Resultado} & \textbf{Fuente}\\
\midrule
Estructura & 8 tablas y 0 vistas & Esquema \codigo{salchipaperia}\\
Productos & 38 productos de la carta & Carta Red Pepper Nueva\\
Pedidos de la muestra & 100 pedidos y 186 líneas & Del 24/10 al 22/11/2026, hora de Lima\\
Importe de la muestra & S/ 2,303.50 & Cantidad por precio aplicado\\
\bottomrule
\end{tabularx}

\seccion{2.}{Relaciones y claves}
Un cliente puede realizar varios pedidos. Cada pedido puede tener varias líneas. El detalle resuelve la relación muchos a muchos entre pedidos y productos con la clave compuesta \codigo{(id\_pedido, id\_producto)}. Una categoría agrupa productos y cada producto pertenece a una categoría.

\begin{tabularx}{\linewidth}{@{}lYY@{}}
\toprule
\textbf{Tabla} & \textbf{Clave y vínculo} & \textbf{Contenido}\\
\midrule
\codigo{clientes} & PK \codigo{id\_cliente} & Nombres y apellidos.\\
\codigo{categorias} & PK \codigo{id\_categoria} & Nombre único y descripción opcional.\\
\codigo{productos} & PK \codigo{id\_producto}; FK categoría & Nombre, precio vigente y disponibilidad.\\
\codigo{snacks} & PK y FK \codigo{id\_producto} & Porción y cremas si se conocen.\\
\codigo{bebidas} & PK y FK \codigo{id\_producto} & Tipo y envase si se conocen.\\
\codigo{hamburguesas} & PK y FK \codigo{id\_producto} & Carne de res o pollo y acompañamiento de papas.\\
\codigo{pedidos} & PK \codigo{id\_pedido}; FK cliente & Fecha, modalidad y estado.\\
\codigo{detalle\_pedido} & PK y FK pedido + producto & Cantidad y precio unitario aplicado.\\
\bottomrule
\end{tabularx}

\nota{PK es la clave primaria. FK es la clave foránea. Las tres tablas de subtipo comparten el identificador del producto.}

\newpage

\seccion{3.}{Normalización e integridad}

\textbf{1FN.} Cada celda registra un valor. Los productos de un pedido se representan en filas del detalle y cada par pedido-producto tiene una fila con la cantidad acumulada.

\textbf{2FN.} La cantidad y el precio aplicado dependen del par pedido-producto. El nombre y el precio vigente del producto dependen del producto y se almacenan en su propia tabla.

\textbf{3FN.} Los atributos de categoría dependen de su identificador. El pedido almacena la clave del cliente y el nombre se obtiene mediante un \codigo{JOIN}. Los importes se calculan a partir de las líneas.

\textbf{FNBC.} Las dependencias declaradas usan claves candidatas como determinantes. En productos también es única la pareja categoría y nombre. La clave compuesta del detalle identifica la cantidad y el precio acordado para ese pedido.

\begin{tabularx}{\linewidth}{@{}lY@{}}
\toprule
\textbf{Regla} & \textbf{Implementación}\\
\midrule
Subtipo del producto & Los disparadores diferidos exigen exactamente una fila en snacks, bebidas o hamburguesas al confirmar la transacción.\\
Bebidas & \codigo{tipo\_bebida} admite Gaseosa o Chicha. Agua y maracuyá conservan su producto y llevan un tipo vacío.\\
Hamburguesas & Las ocho variantes de la carta conservan su clave de producto, su precio, su carne y el acompañamiento de papas.\\
Precio de la línea & \codigo{agregar\_item} toma el precio vigente al crear la línea. Al aumentar unidades, conserva el precio aplicado.\\
Total del pedido & Se obtiene sumando cantidad por precio unitario de cada línea. El resultado depende del detalle y se consulta en SQL.\\
\bottomrule
\end{tabularx}

\seccion{4.}{Migración del esquema}
La migración \codigo{10\_simplificar\_pedidos.sql} retira teléfonos, tipos y series de comprobante, comprobantes, métodos de pago y pagos. También retira las tres vistas anteriores, los campos de documento del cliente, la clasificación interna de categorías, el volumen y la temperatura de las bebidas y el tipo de papa de los snacks.

La migración se ejecuta como una transacción. La prueba de restauración previa recuperó 102 pedidos, 123 filas de pagos y las tres vistas del respaldo. La copia restaurada sirvió para comprobar la migración y los 100 pedidos de la muestra conservaron 186 líneas y S/ 2,303.50. El esquema activo conserva además los dos pedidos iniciales.

\seccion{5.}{Práctica y verificación}
En DataGrip, sincroniza \codigo{salchipaperia} y ejecuta por bloques \codigo{07\_validar\_datagrip.sql}. Las consultas muestran las ocho tablas, clasifican los productos y calculan el total desde \codigo{detalle\_pedido}.

El \href{https://nexlabsgroup.com/basededatos/}{portal Red Pepper} permite seleccionar o crear una comanda, editar su cliente o modalidad, agregar y mover líneas y comparar los valores de antes y después. La sección Resumen del pedido muestra el total. Los ejercicios se guardan en el navegador. El Excel documenta los 100 pedidos, el detalle, los clientes ficticios y los 38 productos.

\end{document}
